\subsubsection{Kosaraju SCC}

\paragraph{Idea intuitiva.}
Encuentra las \textbf{componentes fuertemente conexas} (SCCs) en $O(V + E)$. Algoritmo en 3 pasos:
\begin{itemize}\setlength\itemsep{0pt}
	\item DFS en $G$ rellenando una pila en \textbf{post-orden} (los nodos que terminan tarde quedan arriba).
	\item Construir $G^T$ (grafo transpuesto).
	\item DFS en $G^T$ procesando nodos en orden de la pila (de mayor a menor tiempo); cada DFS da una SCC.
\end{itemize}
La demostracion: dos nodos $u, v$ estan en la misma SCC sii hay camino $u \leadsto v$ y $v \leadsto u$; Kosaraju captura exactamente esto al explorar $G^T$ desde los nodos con mayor tiempo de finalizacion en $G$.

\paragraph{Propiedades clave (invariantes).}
\begin{itemize}\setlength\itemsep{0pt}
	\item Cada SCC es maximal fuertemente conexa.
	\item El \textbf{grafo de SCCs} (metagraf) es un DAG.
	\item Kosaraju y Tarjan dan el mismo resultado; Tarjan es ``mas rapido'' en practica (un solo DFS).
	\item Cada nodo pertenece a exactamente una SCC.
\end{itemize}

\paragraph{Codigo clave.}
La estructura de Kosaraju:
\begin{verbatim}
vector<int> order;
vector<bool> used(n);
function<void(int)> dfs1 = [&](int v) {
    used[v] = true;
    for (int u : g[v]) if (!used[u]) dfs1(u);
    order.push_back(v);  // post-orden
};
function<void(int, int)> dfs2 = [&](int v, int root) {
    comp[v] = root;
    for (int u : gt[v]) if (comp[u] == -1) dfs2(u, root);
};
// 1: DFS en G para llenar order
// 2: DFS en G^T en orden inverso para asignar SCCs
\end{verbatim}

\paragraph{Complejidad.}
\begin{itemize}\setlength\itemsep{0pt}
	\item $O(V + E)$.
	\item Constante 2 (dos DFS + el grafo transpuesto).
\end{itemize}

\paragraph{Donde tocar para modificar.}
\begin{itemize}\setlength\itemsep{0pt}
	\item \textbf{2-SAT}: SCC sobre grafo de implicaciones.
	\item \textbf{Construir el DAG de SCCs}: aniadir aristas entre SCCs diferentes.
	\item \textbf{Tamano de cada SCC}: aniadir contador durante el segundo DFS.
	\item \textbf{Grafo grande}: usar Tarjan para evitar construir $G^T$.
\end{itemize}

\paragraph{Trampas y errores frecuentes.}
\begin{itemize}\setlength\itemsep{0pt}
	\item El snippet usa \texttt{vector<...> adj[n]} (VLA); reemplazable por \texttt{vector<vector<pair<int,int>>} para portabilidad.
	\item El orden del segundo DFS debe ser \textbf{inverso} al de finalizacion del primero.
	\item Si el grafo no es conexo, aniadir un loop sobre todos los nodos en el primer DFS.
	\item Confundir SCC con componentes conexas (las primeras son para dirigidos, las segundas para no dirigidos).
\end{itemize}

\paragraph{Codigo.}
Ver \texttt{cpp.tex}, seccion \textit{Kosaraju SCC}.

\vspace{0.5em}
